\section{Worked Cases II: Mass Category and Assembly}
\label{sec:assembly-cases}

The cases in this section begin after, or explicitly test, the calendar decision.  They show why a permitted requested Mass must be classified before its texts are placed into the Order.

\begin{casefile}{19}{An ordinary IV-class feria repeats the preceding Sunday}
\item[Calendar facts] A weekday outside the ferias with their own Mass is a IV-class feria; no feast or Office of Our Lady on Saturday intervenes.
\item[Ranks and rows] The day is position 28.  The reference to the preceding Sunday comes from rubric 299, not from occurrence.
\item[Impeded item] Nothing is impeded; the feria itself has no separately printed formulary.
\item[Mass category] This is the Mass of the ferial Office assembled from the preceding Sunday's page, not a Sunday Mass celebrated again.
\item[Assembly] Take the preceding Sunday's Introit, Collect, Epistle, intervening chants, Gospel, Offertory, Secret, Communion, and Postcommunion.  Recalculate the conditions for the actual feria: omit Gloria because its Office lacks \latin{Te Deum}; omit Creed because it is not Sunday; use the actual seasonal or Common Preface; add any independently admitted commemorations; use the day's color; and recalculate the dismissal, blessing, and Last Gospel from the actual ferial and seasonal conditions.
\item[Boundary] Do not copy the Sunday worksheet mechanically.  Its proper texts are repeated, but its Sunday status, commemoration ceiling, Gloria, and Creed are not.
\item[Rubric locus] 26; precedence table 91 position 28; 299; 431--432, 475--499, 505, 509.
\end{casefile}

\begin{casefile}{20}{A resumed Sunday after Epiphany placed after Pentecost}
\item[Calendar facts] Septuagesima displaced one of the Third through Sixth Sundays after Epiphany, and rubric 18 assigns it between the Twenty-third and final Sunday after Pentecost.
\item[Ranks and rows] It is a II-class Sunday at position 15; the Twenty-fourth and Last Sunday remains the final Sunday.
\item[Impeded item] This is a rubricated resumption, not transfer of a previously celebrated Sunday and not a choice made because a parish prefers its Gospel.
\item[Mass category] Celebrate the resumed Sunday as the Mass of the Office.
\item[Assembly] From the Twenty-third Sunday after Pentecost take Introit, Gradual, Alleluia with verse, Offertory, and Communion.  From the resumed Epiphany Sunday take Collect, Epistle, Gospel, Secret, and Postcommunion.  Say Gloria and Creed, use the governing Sunday Preface, add only admitted commemorations, and retain the ordinary Sunday conclusion.  Do not take the Twenty-third Sunday's prayers or readings.
\item[Boundary] The actual yearly list of resumptions must be computed from Easter and checked against the Ordo; this case supplies the construction after that computation.
\item[Rubric locus] 18; precedence table 91 position 15; 298; 431, 475, 482--499.
\end{casefile}

\begin{casefile}{21}{A qualified First Friday on a III-class Lenten feria}
\item[Calendar facts] The first Friday is an ordinary III-class Lenten feria, not an Ember day.  In this church or oratory the exercises specified in rubric 385b in honor of the Sacred Heart are actually held that day.
\item[Ranks and rows] The calendar day is position 22; the requested Sacred Heart Mass is a votive of III class under rubrics 384--385, a Mass class rather than a precedence-table row.
\item[Impeded item] The Lenten feria remains the Office and becomes a privileged commemoration in the votive Mass.
\item[Mass category] One of the church's maximum of two permitted Sacred Heart Masses may be celebrated as a III-class votive.
\item[Assembly] Use the white Sacred Heart votive formulary, Gloria, no Creed, the Sacred Heart Preface, and solemn tone if sung.  Add the Lenten feria's Collect, Secret, and Postcommunion first as a privileged commemoration.  Do not add the ferial readings, Tract, color, or prayer over the people: rubric 506 attaches that prayer to Masses \emph{of} Lenten and Passiontide ferias, not to every Mass celebrated in Lent.
\item[Boundary] Private First-Friday observance without the specified exercises does not establish this III-class permission.
\item[Rubric locus] 109e, 113; 384--386; 431--432, 476d, 491, 506.
\end{casefile}

\begin{casefile}{22}{First Friday requests on I-, II-, III-, and IV-class days}
\item[Calendar facts] In each branch the first Friday and required exercises are present; only the day's class changes.
\item[Ranks and rows] A III-class votive is admitted only on liturgical days III and IV.  A I- or II-class day is not converted by the devotion.
\item[Impeded item] On a I- or II-class day celebrate the Office and make no commemoration of the prohibited III-class votive.  On a III- or IV-class day the qualified votive may be selected and the Office is commemorated only as the general and identity rules admit.
\item[Mass category] I: if the day is Sacred Heart itself, use its I-class Office Mass; otherwise use the winning Office.  II: use the Office.  III and IV: up to two Sacred Heart III-class votives in that church.  Without exercises on a IV-class day, only an independently justified IV-class votive may be considered.
\item[Assembly] Office branches follow their own complete proper and receive nothing from the request.  Qualified III-class branches use white Sacred Heart propers, Gloria, no Creed, Sacred Heart Preface, and up to two admitted commemorations.  An unqualified IV-class Sacred Heart votive requires a just cause, omits Gloria and Creed, follows the IV-class prayer limit, and uses ferial tone if sung.
\item[Boundary] The number “two” limits qualifying Sacred Heart Masses in the place; it is not a requirement to schedule two or a faculty for every celebrant.
\item[Rubric locus] 317--319; 384--389; 431--432, 475--476, 491.
\end{casefile}

\begin{casefile}{23}{First Saturday on an otherwise free Saturday}
\item[Calendar facts] The alternative calendar day is a IV-class Saturday feria, so rubric 78 establishes the IV-class Office of Our Lady on Saturday.  The Immaculate-Heart exercises required by rubric 385c are also held.
\item[Ranks and rows] Our Lady on Saturday is position 27.  The requested Immaculate Heart Mass is separately a III-class votive permission.
\item[Impeded item] There are two lawful but distinct paths.  Because both are Marian, whichever is not selected is not commemorated.
\item[Mass category] Path A is the Mass of the BVM-Saturday Office.  Path B is the single qualified Immaculate-Heart III-class votive allowed in that church or oratory.
\item[Assembly] Path A uses the seasonally appointed BVM-Saturday formulary, white, Gloria, no Creed, and the Marian Preface with the appropriate clause.  Path B uses the Immaculate Heart formulary, white, Gloria, no Creed, its governing Marian Preface, and the III-class commemoration limits.  In neither path splice the other Marian Collect, readings, or chants.
\item[Boundary] The first-Saturday devotion does not rename Our Lady on Saturday, and the BVM-Saturday Office does not by itself prove that the special exercises occurred.
\item[Rubric locus] 78; precedence table 91 position 27; 317, 385c--386, 432b, 476d, 495.
\end{casefile}

\begin{casefile}{24}{First Saturday requests on I-, II-, III-, and IV-class days}
\item[Calendar facts] The prescribed Immaculate-Heart exercises occur on first Saturday.  Compare a I-class Marian feast, a II-class non-Marian day, a III-class saint, and the IV-class BVM-Saturday day.
\item[Ranks and rows] The votive is III class, so I and II days prohibit it.  III and IV days admit it subject to rubric 317 and the commemoration rules.
\item[Impeded item] On I and II days nothing of the prohibited votive is commemorated.  On a III non-Marian saint's day, the saint may supply an ordinary commemoration within the two-commemoration ceiling.  On a III Marian day or the IV Marian Saturday Office, select either the Office Mass or votive and exclude commemoration of the other.
\item[Mass category] The I and II branches use the Office; the III and IV branches may use the one qualified III-class votive.
\item[Assembly] Office branches use their complete formularies.  The permitted votive uses the Immaculate Heart proper, white, Gloria, no Creed, its Marian Preface, and only admitted prayer triads.  It never imports the losing Office's readings or chants.
\item[Boundary] Apparition-associated devotional requests explain the gathering but do not enlarge the 1962 permission or overcome a I- or II-class day.
\item[Rubric locus] 317--319, 384--386; 431--432, 476d, 495.
\end{casefile}

\begin{casefile}{25}{First Thursday or First Saturday and Christ the Eternal High Priest}
\item[Calendar facts] Special exercises for sanctification of the clergy are actually held on the first Thursday or first Saturday in the church or oratory.
\item[Ranks and rows] Rubric 385a permits one III-class votive of Jesus Christ, Supreme and Eternal Priest, only on a III- or IV-class day.
\item[Impeded item] On I- or II-class days the request disappears without commemoration.  On a III- or IV-class Office of another mystery of the same Divine Person, rubric 317 permits either the Office or the votive and excludes commemoration of the other.
\item[Mass category] Where eligible, exactly one such III-class votive is allowed in the place.
\item[Assembly] Use the appointed Christ-the-Priest votive formulary and color, Gloria, no Creed, its proper Preface if assigned or otherwise the rubricically governing Preface, and at most two admitted commemorations.  Use solemn tone if sung.  Do not assume a Thursday Eucharistic devotion satisfies the stated clergy-sanctification exercises.
\item[Boundary] The first-Thursday and first-Saturday alternatives share one kind of permission; neither creates a weekly Thursday votive privilege.
\item[Rubric locus] 317--319, 384--386; 431--432, 475--499.
\end{casefile}

\begin{casefile}{26}{An external solemnity on Sunday}
\item[Calendar facts] The Sacred Heart feast occurred during the week and its external solemnity is lawfully kept on an adjacent ordinary II-class Sunday under rubrics 356--360.
\item[Ranks and rows] The Sunday remains the occurring Office at position 15.  The external celebration is a II-class votive Mass, not a transferred calendar feast; a I-class Sunday would prohibit this II-class-votive branch and require a fresh calculation.
\item[Impeded item] Nothing is transferred.  The Sunday Office remains intact.  Because a Sunday and a mystery of the Lord exclude one another's commemoration, no Sunday prayer triad is appended to the Sacred Heart votive.
\item[Mass category] The place may celebrate one sung and one Low Mass, or two Low Masses, of the external solemnity; other Masses follow their independently applicable rules.
\item[Assembly] For an authorized external-solemnity Mass use the white Sacred Heart formulary, Gloria, Sacred Heart Preface, and no Creed by votive class alone; add the Creed because a Sunday occurs.  Apply the II-class-votive prayer limit and ordinary Order.  The Office, Breviary, and other Sunday Masses remain those of Sunday.
\item[Boundary] Only feasts listed in rubric 358 or covered by a produced indult possess this faculty.  Parish custom cannot manufacture an external solemnity.
\item[Rubric locus] 112b; 343; 356--361; 431--432, 475a, 476c, 491.
\end{casefile}

\begin{casefile}{27}{A funeral on an ordinary Sunday and on an Advent Sunday}
\item[Calendar facts] One Requiem directly connected with the funeral is proposed.  Compare an ordinary II-class Sunday with a I-class Advent Sunday, assuming in the ordinary-Sunday branch that no rubric-393 bar, rubric-407 external solemnity, or rubric-406 named local or obligatory feast applies.
\item[Ranks and rows] The funeral Mass is a I-class Requiem category.  Its admission is nevertheless governed first by rubric 393 and then by rubrics 406--407: positions 1--6, the other named local or obligatory feasts, and a Sunday external solemnity prohibit it.
\item[Impeded item] The Sunday Office remains the day's Office in both branches.  On the qualified ordinary II-class Sunday, the 1960 text does not prohibit the funeral Requiem outside the order of that Office merely because it is Sunday, and the Office is not commemorated in the Requiem.  On Advent Sunday, position 6 prohibits the funeral Mass; that Mass may be transferred to the nearest similarly unimpeded day under rubric 408.
\item[Mass category] Admitted branch: the one funeral Requiem.  Prohibited branch: the Advent Sunday Mass of the Office at the funeral action, with any later Requiem separately authorized and recalculated.
\item[Assembly] Admitted Requiem: black, Requiem formulary for the deceased, one proper prayer, \latin{Dies irae}, no Gloria or Creed, Requiem Preface and Agnus Dei, \latin{Requiescant in pace}, no blessing, and no Last Gospel if absolution at the bier follows; make no Sunday commemoration.  Prohibited branch: use the full violet Advent Sunday proper, omit Gloria, say Creed because it is Sunday, and make no Requiem commemoration.
\item[Boundary] The ordinary-Sunday result is edition-textual, conditional on the absence of every rubric-393, 406, and 407 bar, and pastorally sensitive.  Actual celebration requires the competent Ordo, present law, and authority; this manual does not grant permission.
\item[Rubric locus] 390--401, 402, 405--409, especially 393 and 406--407; precedence table 91 positions 1--6 and 15; 432a, 432d, 475a, 476f, 499, 507c--510.
\end{casefile}

\begin{casefile}{28}{Death-day, anniversary, and daily Requiems}
\item[Calendar facts] Compare a Mass on the day of death, an annual anniversary, and an otherwise requested daily Requiem.
\item[Ranks and rows] These are respectively Requiems of II, III, and IV class.  Their class is a Mass-admission category, not the class of the calendar day.
\item[Impeded item] The II-class death-day Mass is prohibited by every Sunday or I-class day.  The III-class anniversary is prohibited on I- and II-class days and can move to the nearest similarly unimpeded day under its rule.  The IV-class daily Requiem is allowed only on IV-class ferias outside Christmas time.
\item[Mass category] Select the exact Requiem kind and its application requirement; never substitute the more permissive funeral rule.
\item[Assembly] Use the formulary assigned to the kind of deceased and occasion, the proper Requiem prayer, no Office commemoration, no Gloria or Creed, optional \latin{Dies irae} in classes II--IV, Requiem Preface and variants, \latin{Requiescant in pace}, no blessing, and the Last Gospel unless a listed omission applies.
\item[Boundary] A “month's mind,” foundation anniversary, cemetery Mass, and Mass within the eight-day period counted inclusively from All Souls each has its own definition under rubrics 415--422.
\item[Rubric locus] 390--401, 410--423, especially 411, 416--419, 423.
\end{casefile}

\begin{casefile}{29}{A priest celebrates one, two, or three Masses on All Souls}
\item[Calendar facts] All Souls lawfully occurs on its proper seat and no question of Sunday transfer remains.
\item[Ranks and rows] It is the I-class Requiem day at precedence position 8; its three-Mass faculty is specific to the day.
\item[Impeded item] No Office commemoration is made.  A funeral Mass on the same day uses the formulary adjustment in rubric 409.
\item[Mass category] One Mass uses the first formulary; two use first then second; three use all three.  A sung or conventual Mass uses the first, with rubric 404 governing order when several are sung or celebrated in different churches.
\item[Assembly] Each Mass uses its own printed chants, readings, and prayer triad; no Gloria, Creed, blessing, or Office commemoration; Requiem Preface and variants.  When three are said without interruption, \latin{Dies irae} is obligatory only in the principal, otherwise the first, and may be omitted in the others unless sung.  Follow the edition's ablution and celebration directions rather than treating three formularies as one long Mass.
\item[Boundary] The faculty does not authorize three Requiems on another date or three applications freely assigned contrary to the governing law.
\item[Rubric locus] 96b, 390--404, 409; 432d, 476f, 499, 507c--510.
\end{casefile}

\begin{casefile}{30}{A wedding on Sunday, All Souls, or during the Triduum}
\item[Calendar facts] The marriage and solemn nuptial blessing are otherwise lawful.  Compare an ordinary Sunday with All Souls and the Sacred Triduum.
\item[Ranks and rows] The Nuptial Mass is a II-class votive but is expressly prohibited on Sundays.  All Souls and the Triduum prohibit the votive, its commemoration, and the nuptial blessing within Mass.
\item[Impeded item] On Sunday, celebrate the occurring Mass, join the Nuptial Collect to its Collect under one conclusion, and retain the appointed nuptial blessings.  On All Souls or in the Triduum, omit all three; the Mass with blessing may be transferred after the marriage as rubric 380 directs.
\item[Mass category] Sunday branch: Mass of the Office with permitted ritual additions.  All Souls/Triduum branch: the unique occurring liturgy only.
\item[Assembly] On Sunday use its complete proper, Gloria, Creed, color, and Preface; join the Nuptial Collect to the occurring Collect under one conclusion, but do not add the Nuptial Secret or Postcommunion, which rubric 380 does not prescribe in this impeded branch.  Pray \latin{Propitiare} and \latin{Deus, qui potestate} after the Pater before \latin{Libera nos}; after receiving the Precious Blood, communicate the spouses; and say \latin{Deus Abraham} after \latin{Ite, missa est} before the general blessing.  Use none of these insertions in the prohibited branch.
\item[Boundary] Presence of the couple does not create a Nuptial Mass permission.  The blessing's personal conditions and any closed-time permission must also be checked.
\item[Rubric locus] 341, 343, 378--381; 444--446; printed \latin{Pro sponsis} formulary.
\end{casefile}

\begin{casefile}{31}{An Ember Saturday with five lessons}
\item[Calendar facts] The worksheet distinguishes three branches: an Ember Saturday whose feria supplies the Office; an ordination on an Ember Saturday when a I- or II-class feast supplies the Office under rubric 300; and another sung or Low Mass.
\item[Ranks and rows] Advent, Lenten, and September Ember Saturdays are II-class ferias at position 18 under rubric 24b; Pentecost Ember Saturday belongs to the I-class octave structure under rubric 66.  An occurring I- or II-class feast can therefore create rubric 300's exceptional ordination branch.
\item[Impeded item] In the ordinary branch, admitted commemorations affect prayer placement, not the lesson count.  In the exceptional branch, celebrate the Ember-Saturday ordination Mass even though it is not the Mass of the feast's Office, and make the feast's admitted commemoration.
\item[Mass category] When the Ember feria supplies the Office, its Mass is the Mass of the Office.  When a I- or II-class feast supplies the Office, rubric 300 nevertheless mandates the exceptional Ember-Saturday Mass for the ordination; that Mass must not be mislabeled a Mass of the feast's Office.
\item[Assembly] In a conventual or ordination Mass, retain all five lessons before the Epistle with their prayers and verses.  In the full form, place the admitted commemorations and the prayer for conferral of sacred Orders after the prayer before the final lesson or Epistle; join the ordination prayer under one conclusion as rubrics 442, 444, and 447b direct, and under rubric 448 admit no other prayers except privileged commemorations.  In another Mass, rubric 468 permits the abbreviated path: first Office prayer with its genuflection when prescribed, first lesson and verses, then the second prayer without \latin{Flectamus genua}, admitted commemorations, and the final lesson or Epistle with its following chant.  Continue with the Gospel and remaining proper; the Ember formulary and its actual seasonal branch, not the commemorated feast, govern color, Gloria, Creed, Preface, and dismissal.
\item[Boundary] A hand missal's abbreviated layout must not erase the full form where the rubric makes it obligatory, and the exceptional ordination rule must not be generalized to a non-ordination Mass.
\item[Rubric locus] 24b, 66, 300; precedence table 91 positions 10 and 18; 442, 444, 447--448, 467--468.
\end{casefile}

\begin{casefile}{32}{A temporal Passiontide Mass versus a saint's Mass in Passiontide}
\item[Calendar facts] Compare the temporal Mass of an ordinary Passiontide feria with a superior saint's feast lawfully celebrated during the same season.
\item[Ranks and rows] The calendar calculation has already selected the feria in one branch and the saint in the other; seasonal location alone does not make both Masses “of the season.”
\item[Impeded item] In the saint branch the Passiontide feria is privilegedly commemorated where admitted, but it does not convert the feast Mass into a temporal Mass.
\item[Mass category] Branch A is a Mass of the temporal Office.  Branch B is a feast's Mass of the Office.
\item[Assembly] In Branch A omit Psalm 42 and the Introit \latin{Gloria Patri} under rubrics 425a and 428, and use the temporal proper.  In Branch B retain Psalm 42 and the Introit doxology because those omissions are stated for \emph{Missae de Tempore}; use the saint's proper/Common and add only the ferial prayer triad.  If the saint lacks a proper Preface, the Holy Cross Preface still governs seasonally under rubric 487b.
\item[Boundary] This distinction does not alter the separate full-preparation omissions after Palm, ashes, candle, Rogation, or other processions under rubric 424.
\item[Rubric locus] 108--109e; 424--429; 482, 487.
\end{casefile}

\begin{casefile}{33}{Two overlays: prayer over the people and proper octave Canon forms}
\item[Calendar facts] First consider a Lenten ferial Mass with two commemorations.  Separately consider a saint or permitted votive Mass celebrated within the Christmas, Easter, or Pentecost octave.
\item[Ranks and rows] Each underlying Mass has already been admitted.  The overlays come from the day or season and do not replace its proper.
\item[Impeded item] Nothing further is displaced.  The issue is whether the stable Order receives a required addition.
\item[Mass category] The first is a Lenten feria's Mass of the Office; the second remains the selected feast or votive.
\item[Assembly] Lenten branch: after the Mass Postcommunion and two commemoration Postcommunions, add \latin{Oremus. Humiliate capita vestra Deo} and the printed prayer over the people under its own conclusion; it remains required even after three Postcommunions.  Octave branch: use the selected Mass's chants, readings, prayers, Preface, and color.  During the Christmas octave insert its proper \latin{Communicantes}; during the Easter and Pentecost octaves insert both the proper \latin{Communicantes} and the proper \latin{Hanc igitur}.  These forms apply even if the selected Mass is not of the octave and has its own Preface.
\item[Boundary] Do not generalize either overlay: the prayer belongs to Lenten and Passiontide ferial Masses outside the Triduum; Christmas has no proper octave \latin{Hanc igitur}, while Easter and Pentecost have both named forms.
\item[Rubric locus] 64--70; 433--446, 482--501, especially 501, 505--506.
\end{casefile}

\begin{casefile}{34}{A saint's page supplies only part of the formulary}
\item[Calendar facts] A saint's Office has won, but the Proper of Saints prints only some prayers and directs the reader to a named Common.
\item[Ranks and rows] The saint's universal or proper row and class have already been established from the calendar; the Common does not confer rank.
\item[Impeded item] Any losing day contributes only an admitted commemoration, never material chosen to fill gaps in the feast.
\item[Mass category] This is the feast's Mass of the Office.
\item[Assembly] Make a source ledger.  Copy each proper item exactly; take every absent Introit, reading, chant, Gospel, Offertory, or Communion from the directed Common; use an allowed alternative only under rubric 305; then add the feast's prayer triad, admitted commemorations in order, color, Gloria and Creed decisions, and Preface under 482.  Apply seasonal Alleluia or Tract directions and retain the ordinary Canon unless the Missal prints an intervention.
\item[Boundary] Do not use a biographical guess to select another Common or treat the Common's most familiar reading as if it were proper to the saint.
\item[Rubric locus] 269--273, 301--305; 427--510.
\end{casefile}

\begin{casefile}{35}{A church with only one Mass}
\item[Calendar facts] A requested Requiem or votive appears generally admissible on the day's class, but this church has only one Mass and also has a conventual obligation, candle or ash blessing, or required Rogation Mass.
\item[Ranks and rows] The calendar day and requested Mass category are calculated separately; the one-Mass constraint is a further admission rule.
\item[Impeded item] The requested Mass is prohibited when the sole Mass must satisfy the conventual obligation and no other priest can do so, or on the named days when the sole Mass must follow the blessing or Litanies.  The request does not survive as a commemoration unless its own rubric expressly grants one.
\item[Mass category] Celebrate the required conventual or ritual Mass.  A different result is possible only if another Mass lawfully fulfills the obligation.
\item[Assembly] Use the complete proper of the required Office or Rogation action and follow the preceding blessing or procession's omission rules at the altar.  Add only commemorations admitted for that Mass category.  Do not retain Requiem color, prayers, or conclusion when the Requiem is prohibited.
\item[Boundary] “The day is IV class” proves too little.  Number of Masses, choral duty, exposition, and an inseparable public action can still bar an otherwise eligible request.
\item[Rubric locus] 289--304, 326; 346--355; 393, especially 393b--c; 424.
\end{casefile}
